\section{Geometry of an isolated degenerate subspace}
\label{sec:bundle_geometry}

We now fix the geometric objects whose statistics can be compared without
choosing a basis in the protected multiplet.  Let $\mathcal H$ be a finite
Hilbert space and let $H(\boldsymbol\lambda)$ be smooth on a coordinate patch
of $\mathcal M$.  We assume that Eq.~\eqref{eq:protected_projector} holds with
constant rank $D$ and that
\begin{equation}
 \Delta(\boldsymbol\lambda)
 =\operatorname{dist}\!\left(
 E_0(\boldsymbol\lambda),
 \operatorname{spec}\bigl[QH(\boldsymbol\lambda)Q\bigr]
 \right)>0,
 \qquad Q=1-P,
 \label{eq:external_gap}
\end{equation}
throughout the patch.  The fibers $P(\boldsymbol\lambda)\mathcal H$ form a
rank-$D$ complex vector bundle.  A local frame $\Psi$ obeys
$\Psi^\dagger\Psi=I_D$ and $P=\Psi\Psi^\dagger$.  On an overlap of patches it
transforms as
\begin{equation}
 \Psi(\boldsymbol\lambda)\longmapsto
 \Psi(\boldsymbol\lambda)U(\boldsymbol\lambda),
 \qquad U(\boldsymbol\lambda)\in U(D).
 \label{eq:fiber_gauge}
\end{equation}
This is a change of coordinates in the same physical subspace.

The Wilczek--Zee connection and curvature in this convention are
\begin{align}
 A_\mu&=i\Psi^\dagger\partial_\mu\Psi,\nonumber\\
 F_{\mu\nu}
 &=\partial_\mu A_\nu-\partial_\nu A_\mu-i[A_\mu,A_\nu],
 \label{eq:connection_curvature}
\end{align}
which transform as a connection and by conjugation, respectively.  Their
Abelian rank-one antecedent and non-Abelian extension are standard
consequences of adiabatic transport~\cite{berry1984,wilczekzee1984,kato1950}.
Connection entries at one point are gauge dependent; curvature eigenvalues,
invariant traces, and Wilson-loop conjugacy classes are not.

The response leaving the fiber is
\begin{equation}
 X_\mu=Q\partial_\mu\Psi
 \in\operatorname{Hom}(P\mathcal H,Q\mathcal H).
 \label{eq:offfiber_response}
\end{equation}
It transforms as $X_\mu\mapsto X_\mu U$.  Differentiating the eigenvalue
equation and projecting with $Q$ gives the reduced-resolvent form
\begin{equation}
 X_\mu=-R_Q(\partial_\mu H)\Psi,
 \qquad
 R_Q=Q(H-E_0)^{-1}Q .
 \label{eq:kato_response}
\end{equation}
The inverse in Eq.~\eqref{eq:kato_response} is restricted to $Q\mathcal H$.
Thus a closing external gap is a singularity of the construction rather than
an unusually large sample from a fixed response ensemble.

The matrix-valued quantum geometric tensor (QGT) is the fiber endomorphism
\begin{equation}
 \mathcal Q_{\mu\nu}
 =(\partial_\mu\Psi)^\dagger Q(\partial_\nu\Psi)
 =X_\mu^\dagger X_\nu .
 \label{eq:matrix_qgt}
\end{equation}
Its Hermitian symmetric and antisymmetric parts give the quantum metric and
Berry curvature,
\begin{align}
 G_{\mu\nu}
 &=\frac{1}{2}\left(
 \mathcal Q_{\mu\nu}+\mathcal Q_{\nu\mu}
 \right),\nonumber\\
 F_{\mu\nu}
 &=i\left(
 \mathcal Q_{\mu\nu}-\mathcal Q_{\nu\mu}
 \right).
 \label{eq:qgt_parts}
\end{align}
For $D=1$, Eq.~\eqref{eq:qgt_parts} reduces to
$g_{\mu\nu}=\operatorname{Re}\mathcal Q_{\mu\nu}$ and
$F_{\mu\nu}=-2\operatorname{Im}\mathcal Q_{\mu\nu}$, with the stated
connection convention~\cite{provostvallee1980,kolodrubetz2017}.  For
$D>1$, each block transforms as
$\mathcal Q_{\mu\nu}\mapsto U^\dagger\mathcal Q_{\mu\nu}U$.

It is useful to separate the central and traceless curvature,
\begin{align}
 F_{\mu\nu}&=F^{(0)}_{\mu\nu}I_D
 +\widetilde F_{\mu\nu},\nonumber\\
 F^{(0)}_{\mu\nu}&=\frac{1}{D}\operatorname{Tr}F_{\mu\nu},
 \qquad
 \operatorname{Tr}\widetilde F_{\mu\nu}=0.
 \label{eq:trace_traceless}
\end{align}
The trace controls the determinant line bundle and, on a closed two-cycle
$\Sigma\subset\mathcal M$, the first Chern number
\begin{equation}
 C_1[\Sigma]=\frac{1}{2\pi}
 \int_\Sigma\operatorname{Tr}F .
 \label{eq:first_chern}
\end{equation}
This normalization follows the many-body Berry-curvature convention of
Ref.~\cite{niu1985}.
The traceless part records internal $SU(D)$ rotation not fixed by $C_1$.
Keeping these pieces separate prevents a topological mean curvature from
being mistaken for fluctuating non-Abelian geometry.

Local statistical statements must be made from conjugation invariants or
from covariantly compared matrices.  At one point these include
$D^{-1}\operatorname{Tr}\widetilde F^k$, the spectrum of
$\widetilde F$, and invariant contractions of $\mathcal Q$.  At different
points, matrices must first be parallel transported to a common fiber, or the
comparison must be restricted to scalar invariants.  Raw entries of $A_\mu$,
$F_{\mu\nu}$, or $X_\mu$ in an arbitrary numerical frame are not primary
observables.

The geometry also separates protection from mixing.  The rank $D$ and the
gap in Eq.~\eqref{eq:external_gap} state that a multiplet exists.  The span of
the responses $X_\mu$, the algebra generated by $X_\mu^\dagger X_\nu$, and
their higher connected moments state how extensively accessible
deformations move that multiplet.  The first set of facts can be exact while
the second remains structured, reducible, or random.  This separation is the
starting point for the definitions below.
